\section{Introduction}
Recent advancements in 3D Gaussian Splatting (3DGS) have enabled real-time rendering and semantic scene understanding for both static and dynamic environments. Extensions of 3DGS now incorporate features from large-scale 2D vision-language models for semantic reasoning. To enable dynamic scene representations, recent methods~\cite{DNERF,dynerf,dynamic3dgs,4dgs,gs-flow,fully} augment Gaussians with temporal modeling to capture changes in geometry and appearance over time. These developments make 3DGS as a promising backbone for applications in robotics, AR/VR, and digital content creation, where efficiency and semantic accuracy are critical. However, the potential of 3DGS in such real-world settings is limited by the substantial resources required for training, especially when scaling to large-scale environments, densely instrumented multi-camera systems, or long-duration dynamic sequences. The key challenge lies in the vast number of candidate viewpoints: although many are redundant or uninformative, choosing which views are most informative is non-trivial. Unlike static reconstruction, dynamic environments introduce additional uncertainty due to evolving geometry, appearance changes, and semantic inconsistencies across time. 

A natural solution is \textbf{next-best-view} (NBV) selection, which seeks to prioritize views that most improve reconstruction and understanding. Effective NBV selection can significantly improve both geometric reconstruction and semantic consistency, especially in dynamic environments where scene content changes over time. In practice, however, existing strategies often rely on random sampling, uncertainty heuristics~\cite{CSS, activeneural}, or black-box selection models~\cite{dnrselect, activeinitsplat}, which do not explicitly maximize information gain and therefore can lead to suboptimal reconstruction and semantic coverage. Moreover, prior work predominantly focuses on static scenes, leaving dynamic NBV selection largely unexplored. Under dynamic scenes, these strategies struggle to generalize across scenes with complex motion and appearance variation, often lead to wasted computation, incomplete reconstructions, and semantic drift, leaving us without a principled way to allocate training resources effectively. 

In this work, we present a unified framework that combines dynamic and semantic 3DGS and formulates NBV selection as an active learning problem. Building on this formulation, we develop the first NBV algorithm tailored to dynamic semantic 3DGS, jointly capturing geometric, semantic, and temporal deformation information. Unlike FisherRF~\cite{fisherrf}, which considers only geometric information for static scenes, our method systematically leverages Fisher Information to identify views that maximally improve both reconstruction quality and dynamic scene modeling. To handle the computational challenges of large-scale dynamic semantic 3DGS, we propose an efficient diagonal approximation of the Fisher Information matrix and introduce a novel estimator for the Fisher Information of deformation networks using the trace of the gradient outer product. Beyond proposing a new algorithm, our framework integrates NBV selection directly into the 3DGS backbone, greatly extending the scenarios and applications of NBV to semantic and dynamic scene understanding. 

We evaluate our framework on large-scale static datasets (Replica~\cite{replica}) and dynamic multi-camera sequences (Neu3D~\cite{dynerf}), showing consistent improvements over baseline methods. Our contributions include:
\begin{itemize}
    \item The first Fisher Information-driven NBV selection framework for \emph{dynamic semantic 3DGS}, capturing geometric, photometric, semantic, and deformation informativeness.
    \item An efficient Fisher Information formulation for large-scale semantic and dynamic 3D Gaussian Splatting, which extends the diagonal approximation of FisherRF~\cite{fisherrf} to semantic Gaussian parameters and deformation networks, enabling tractable NBV selection beyond static radiance fields.
    \item A novel method to estimate Fisher Information for deformation MLPs, leveraging the trace of the gradient outer product as a proxy.
\end{itemize}
